\section{Exact profile calculus}\label{supp:profile}

For cluster $i$, write
\[
 \ell_{i,h}(\beta,\gamma)
 =m_i^{-1}\sum_{j=1}^{m_i}\rho_{\tau,h}(Y_{ij}-X_{ij}^{\mathsf T}\beta-Z_{ij}^{\mathsf T}\gamma)
 +\tfrac12\gamma^{\mathsf T}\Lambda\gamma.
\]
Let $\widehat\gamma_i(\beta)=\argmin_{\gamma}\ell_{i,h}(\beta,\gamma)$ and
$L_{n,h}^{\mathrm{prof}}(\beta)=n^{-1}\sum_i\ell_{i,h}\{\beta,\widehat\gamma_i(\beta)\}$.

\begin{proposition}[Profile gradient and Hessian]\label{prop:supp-profile}
Let
\[
 W_i=m_i^{-1}\diag\{K_h(r_{ij})\},\qquad
 B_i=Z_i^{\mathsf T}W_iZ_i+\Lambda,\qquad
 A_i=B_i^{-1}Z_i^{\mathsf T}W_iX_i.
\]
Then
\begin{equation}\label{eq:supp-gradient}
 \nabla L_{n,h}^{\mathrm{prof}}(\beta)
 =-n^{-1}\sum_i m_i^{-1}X_i^{\mathsf T}\psi_{\tau,h}(r_i),
\end{equation}
\begin{equation}\label{eq:supp-hessian}
 \nabla^2L_{n,h}^{\mathrm{prof}}(\beta)
 =n^{-1}\sum_i\{X_i^{\mathsf T}W_iX_i-X_i^{\mathsf T}W_iZ_iB_i^{-1}Z_i^{\mathsf T}W_iX_i\},
\end{equation}
and, with $\widetilde X_i=X_i-Z_iA_i$,
\begin{equation}\label{eq:supp-identity}
 X_i^{\mathsf T}W_iX_i-X_i^{\mathsf T}W_iZ_iB_i^{-1}Z_i^{\mathsf T}W_iX_i
 =\widetilde X_i^{\mathsf T}W_i\widetilde X_i+A_i^{\mathsf T}\Lambda A_i.
\end{equation}
\end{proposition}

\begin{proof}
The stationarity equation for the nuisance profile is
\begin{equation}\label{eq:supp-stationarity}
 0=\nabla_{\gamma}\ell_{i,h}\{\beta,\widehat\gamma_i(\beta)\}
 =-m_i^{-1}Z_i^{\mathsf T}\psi_{\tau,h}\{r_i(\beta)\}
 +\Lambda\widehat\gamma_i(\beta).
\end{equation}
Because $B_i\succeq\Lambda\succ0$, the profile map is locally continuously differentiable. The envelope theorem gives
\[
 \frac{d}{d\beta}\ell_{i,h}\{\beta,\widehat\gamma_i(\beta)\}
 =\partial_{\beta}\ell_{i,h}\{\beta,\widehat\gamma_i(\beta)\}
 =-m_i^{-1}X_i^{\mathsf T}\psi_{\tau,h}\{r_i(\beta)\},
\]
which proves \eqref{eq:supp-gradient}.

Differentiate \eqref{eq:supp-stationarity} with respect to $\beta$. The block derivatives of the complete criterion are
\[
 H_{\beta\beta,i}=X_i^{\mathsf T}W_iX_i,
 \qquad H_{\beta\gamma,i}=X_i^{\mathsf T}W_iZ_i,
 \qquad H_{\gamma\gamma,i}=B_i.
\]
Implicit differentiation yields
\[
 \frac{d\widehat\gamma_i(\beta)}{d\beta}
 =-B_i^{-1}Z_i^{\mathsf T}W_iX_i=-A_i.
\]
Differentiating the profile gradient therefore gives the Schur complement in \eqref{eq:supp-hessian}.

For \eqref{eq:supp-identity}, note that $Z_i^{\mathsf T}W_iX_i=B_iA_i$ and $Z_i^{\mathsf T}W_iZ_i=B_i-\Lambda$. Hence
\begin{align*}
 \widetilde X_i^{\mathsf T}W_i\widetilde X_i
 &=X_i^{\mathsf T}W_iX_i-2A_i^{\mathsf T}B_iA_i
 +A_i^{\mathsf T}(B_i-\Lambda)A_i\\
 &=X_i^{\mathsf T}W_iX_i-A_i^{\mathsf T}B_iA_i-A_i^{\mathsf T}\Lambda A_i.
\end{align*}
Adding $A_i^{\mathsf T}\Lambda A_i$ and using
$A_i^{\mathsf T}B_iA_i=X_i^{\mathsf T}W_iZ_iB_i^{-1}Z_i^{\mathsf T}W_iX_i$
proves the result.
\end{proof}

\subsection{Dual profile scores with bounded cluster sizes}

Write $\mathcal C_\tau=[\tau-1,\tau]$ and let
$c_\tau(s)=\sup_v\{sv-\rho_{\tau,1}(v)\}$ for $s\in\mathcal C_\tau$.
The kernel is symmetric and has compact support. Consequently
$\rho_{\tau,h}^*(s)=h c_\tau(s)$ on $\mathcal C_\tau$, and
$c_\tau$ is bounded and continuous on this interval.

\begin{proposition}[Dual profile score]\label{prop:supp-dual-profile}
For $u=Y_i-X_i\beta$ and $h\geq0$, the profiled cluster criterion equals
\begin{equation}\label{eq:supp-profile-dual}
 \max_{s\in\mathcal C_\tau^{m_i}}
 \left\{
 \frac{s^{\mathsf T}u}{m_i}
 -\frac{s^{\mathsf T}Z_i\Lambda^{-1}Z_i^{\mathsf T}s}{2m_i^2}
 -\frac{h}{m_i}\sum_{j=1}^{m_i}c_\tau(s_j)
 \right\}.
\end{equation}
For $h>0$, its maximiser $s_{i,h}$ is the vector
$\psi_{\tau,h}(r_i)$. At $h=0$, every maximiser satisfies
\[
 \Lambda\gamma_i^\dagger=m_i^{-1}Z_i^{\mathsf T}s_{i,0},
 \qquad
 s_{ij,0}\in\partial\rho_\tau
 (Y_{ij}-X_{ij}^{\mathsf T}\beta-Z_{ij}^{\mathsf T}\gamma_i^\dagger).
\]
If $Y_i$ has a joint conditional density given $(X_i,Z_i,m_i)$,
the maximiser $s_{i,0}$ is unique almost surely at each fixed $\beta$.
If additionally
$\E[(m_i^{-1}\sum_j|X_{ij,k}|)^2]<\infty$ for coordinate $k$, then
\[
 m_i^{-1}X_{i,k}^{\mathsf T}s_{i,h}
 \longrightarrow m_i^{-1}X_{i,k}^{\mathsf T}s_{i,0}
 \quad\text{in }L_2\text{ as }h\downarrow0.
\]
The same conclusion holds for a fixed projected score when its design
envelope has a finite second moment.
\end{proposition}

\begin{proof}
Fenchel duality for the finite continuous convex criterion introduces
$s_j\in\mathcal C_\tau$ for each check loss. Minimising over $\gamma$
gives $\gamma=\Lambda^{-1}Z_i^{\mathsf T}s/m_i$ and the quadratic term
in \eqref{eq:supp-profile-dual}. The scaling identity
$\rho_{\tau,h}(v)=h\rho_{\tau,1}(v/h)$ gives its conjugate.
The Fenchel optimality conditions give the stated score and nuisance
stationarity relations.

At $h=0$ the value in \eqref{eq:supp-profile-dual} is a finite convex
Lipschitz function of $u\in\R^{m_i}$. It is differentiable Lebesgue
almost everywhere. Its subdifferential in $u$ consists of the dual
maximisers divided by $m_i$, so differentiability makes the maximiser
unique. The joint conditional density gives this event probability one.
The dual objective converges uniformly on its compact feasible set as
$h\downarrow0$, since $c_\tau$ is bounded. Every accumulation point of
$s_{i,h}$ therefore maximises the limiting objective, giving
$s_{i,h}\to s_{i,0}$ almost surely. The bound
$|s_{ij,h}|\leq1$ and the assumed design envelope permit dominated
convergence for the squared coordinate or projected score.
\end{proof}

When a profiled residual equals zero, its dual coordinate is chosen
jointly with the other residuals through nuisance stationarity. This
selection enters both the population profile gradient and its limiting
variance. Joint outcome densities permit the preceding argument even
when the fitted profile residual has positive probability at zero.

